\section{Miscellaneous remarks}\label{sec:misc}

\subsection{Behaviour of the gauge field near flow time zero}

In the regularized theory, the time-dependent field $B_{\mu}(t,x)$
satisfies the boundary condition \eqref{eq:2.3} and its correlation functions
thus converge to those of the bare gauge field $A_{\mu}(x)$
when $t$ goes to zero. However,
since the latter requires renormalization by a divergent constant,
the correlation functions tend to become singular at $t=0$
after renormalization and removal of the regularization.
Their asymptotic behaviour for $t\to0$ is then
described by an expansion
\begin{equation}
   B_{\mu}(t,x)=c_B(t)(\ren{A})_{\mu}(x)+\rmO(t)
   \label{eq:8.1}
\end{equation}
in local renormalized fields with singular coefficients \cite{Symanzik}.
At one-loop order of perturbation theory, for example,
the coefficient $c_B(t)$ is found to
diverge logarithmically at $t=0$.

It is straightforward to show that the
renormalization group equation
\begin{equation}
   \left\{\mu{\partial\over\partial\mu}+\beta{\partial\over\partial g}
   -2\gamma\lambda{\partial\over\partial\lambda}
   +\gamma+\beta/g
   \right\}c_B(t)=\rmO(t)
   \label{eq:8.2}
\end{equation}
holds, where
\begin{equation}
   \beta=-b_0g^3+\rmO(g^5),
   \qquad
   \gamma=c_0g^2+\rmO(g^4),
   \label{eq:8.3}
\end{equation}
are the beta function and the anomalous dimension of the
gauge field. In the Landau gauge, for example,
the equation can be easily integrated and one finds that
\begin{equation}
   B_{\mu}(t,x)
   \mathrel{\mathop\sim_{t\to0}}
   \bigl(2b_0\gbar(q)^2\bigr)^{1/2-c_0/2b_0}
   R(g)(\ren{A})_{\mu}(x),
   \label{eq:8.4}
\end{equation}
$\gbar(q)$ being the running coupling at momentum $q=(8t)^{-1/2}$
and $R(g)$ the factor that relates the renormalized to
the renormalization-group-invariant gauge field.

Equation \eqref{eq:8.4} is a remnant of the
boundary condition satisfied by the gauge field in the regularized theory.
It shows that the field generated by the flow equation
is connected to the fundamental field in a universal manner, i.e.~there
is no room for finite renormalizations here, the reason being
that any such renormalization would violate the BRS symmetry.

\subsection{Gauge-invariant composite fields}

Wilson loops and gauge-invariant local fields are independent of
the parameter $\alpha_0$ in the flow equation
and do not require renormalization
at positive flow time. At small flow times,
their asymptotic behaviour is determined by the scaling
properties of the corresponding
renormalized fields in the $\SUn$ gauge theory and thus
reflects the singular nature of the latter.

For illustration, consider the density
\begin{equation}
   E=\frac{1}{4}G_{\mu\nu}^aG_{\mu\nu}^a
   \label{eq:8.5}
\end{equation}
that has previously been studied in ref.~\cite{WilsonFlow}.
Since $E$ can mix with the unit field, the dominant term
at small flow times,
\begin{align}
   E(t,x)=\langle E(t,x)\rangle+c_E(t)
   \ren{\Bigl\{\frac{1}{4}F^a_{\mu\nu}F^a_{\mu\nu}\Bigr\}}(x)+\rmO(t),
   \label{eq:8.6}
\end{align}
is its expectation value, while the first subleading term is proportional
to the renormalized action density of the fundamental field.
The associated coefficients,
$\langle E(t,x)\rangle$ and $c_E(t)$,
satisfy a renormalization
group equation that allows their exact asymptotic behaviour
at small $t$ to be worked out analytically.

\subsection{Theories with matter fields}

In the presence of matter fields, the gradient flow is
defined by the same equations
as in the pure gauge theory (sect.~\ref{ssec:defflow}).
The interesting but rather more complicated case where
the matter fields are included in the time evolution
is not considered here.

In a renormalizable theory, matter fields are
scalar fields of dimension $1$ or fermion fields of
dimension $3/2$. Provided
the regularization preserves the
gauge symmetry, our argumentation
in sects.~\ref{sec:brs} and \ref{sec:fin} then carries over literally. In particular,
additional boundary counterterms involving the matter fields
are excluded by the global symmetries and the requirement
that their dimension must be less than or equal to $4$.
The finiteness of the correlation functions
at positive flow times is therefore again guaranteed.

QCD is a prominent example of a gauge theory that
has all the required properties to ensure finiteness,
but the same applies to many more theories of interest,
including the SU(2) Higgs model,
supersymmetric versions of QCD and technicolour theories.

\subsection{Lattice regularization}

While the Feynman rules are more
complicated than with dimensional regularization,
the perturbative analysis of the gradient flow on the lattice
is not expected to run into fundamental difficulties.

A possible choice of the gradient term
in the lattice flow equation
is the gradient of the Wilson action
\cite{WilsonFlow}.
There is however no reason to choose this particular action
or to require that it coincides with the gauge action
of the theory.
As long as the term has the
correct form and normalization in the classical continuum limit
(and thus at tree-level of perturbation theory),
the correlation functions of the time-dependent gauge field
will not depend on the exact choices one makes, except
for lattice effects vanishing proportionally to a positive power
of the lattice spacing $a$.

The continuous Lorentz symmetry is broken on the lattice,
but the remaining exact symmetries are sufficient to exclude
counterterms that have not already appeared
in the continuum theory.
We therefore expect that the correlation functions
of the time-dependent fields
do not require renormalization and
that their values in the continuum limit
are independent of the regularization
(up to finite renormalizations of
the coupling and the gauge-fixing parameter).
Moreover, if the lattice theory is $\rmO(a)$ improved,
the correlation functions are automatically improved too,
because there are no candidate $\rmO(a)$ counterterms
with all the required properties.
